\documentclass[11pt]{article}
% =========================================================
%  學生講義 共用前言（以 XeLaTeX 編譯）
%  需要套件：xeCJK, tcolorbox（其餘多在 BasicTeX 內）
% =========================================================
\usepackage[a4paper,margin=2.3cm]{geometry}
\usepackage{amsmath,amssymb}
\usepackage{xcolor}
\usepackage{graphicx}

% ---- 中文字型（macOS 系統字型）----
\usepackage{xeCJK}
\setCJKmainfont{Songti TC}      % 內文：宋體
\setCJKsansfont{Heiti TC}       % 標題/強調：黑體
\xeCJKsetup{AutoFakeBold=2.2, AutoFakeSlant=0.2}

% ---- 版面 ----
\setlength{\parindent}{0pt}
\setlength{\parskip}{0.55em}
\linespread{1.15}
\renewcommand{\baselinestretch}{1.15}

% ---- 顏色 ----
\definecolor{cBlue}{HTML}{1f6feb}
\definecolor{cGray}{HTML}{6b7280}
\definecolor{cGrayBg}{HTML}{f3f4f6}
\definecolor{cPurp}{HTML}{7a3ff2}
\definecolor{cPurpBg}{HTML}{f5f1ff}

% ---- 標註框 ----
\usepackage[most]{tcolorbox}

% 就地補充新概念（灰底）：\begin{補充框}{標題}
\newtcolorbox{補充框}[1]{
  enhanced, breakable, arc=2pt, boxrule=0.6pt,
  colback=cGrayBg, colframe=cGray,
  left=9pt,right=9pt,top=7pt,bottom=7pt,
  before upper={\textbf{\sffamily #1}\par\vspace{3pt}},
}

% 延伸 / 開放問題（紫色左邊條）：\begin{延伸框}{標題}
\newtcolorbox{延伸框}[1]{
  enhanced, breakable, arc=2pt, boxrule=0pt,
  colback=cPurpBg, colframe=cPurpBg,
  borderline west={3pt}{0pt}{cPurp},
  left=10pt,right=9pt,top=7pt,bottom=7pt,
  before upper={\textbf{\sffamily 延伸閱讀｜#1}\par\vspace{3pt}},
}

% ---- 圖：置中、底下小字說明 ----
% \fig{檔名}{寬度}{說明}
\newcommand{\fig}[3]{%
  \begin{center}
  \includegraphics[width=#2\linewidth]{#1}\\[2pt]
  {\small\color{cGray} #3}
  \end{center}}

% ---- 章節標題用黑體 ----
\newcommand{\h}[1]{\sffamily #1}


\begin{document}

\begin{center}
{\sffamily\Huge\bfseries 數1-4 直線與圓}\\[3pt]
{\sffamily\large 普通高中 數學（必修・第一冊）・學生講義}
\end{center}
\vspace{0.6em}

本章把「圖形」翻譯成「方程式」。一條直線、一個圓本來是純幾何的對象，只要在平面上鋪一層坐標格子，它們就能用代數式表示，並用計算回答幾何問題：兩條直線是否相交、一點離一條直線多遠、一條直線是否與圓相交、相交於幾點。這套「用坐標研究幾何」的方法稱為\textbf{解析幾何（analytic geometry）}。

本章是高中第一次把幾何與代數正式接起來。其中的兩點距離、斜率、距離公式，之後會在三角比中接上傾斜角的正切，並在平面向量中以向量重寫；圓與直線的位置判定，則是日後圓錐曲線位置判定的原型。本章以\textbf{定性理解與基本計算}為主。

\medskip
本章主線：

\begin{center}
坐標平面與對稱 $\rightarrow$ 直線方程式（斜率、三式、平行與垂直、點到直線距離）$\rightarrow$ 圓方程式 $\rightarrow$ 直線與圓的位置關係
\end{center}

\section*{\sffamily 4-1\quad 坐標平面與對稱}

\subsection*{\sffamily A.\ 把幾何鋪上坐標}

平面上一個點本來只是一個位置。\textbf{笛卡兒（Descartes）}的方法是：取一個原點與兩條互相垂直的數線（$x$ 軸、$y$ 軸），則每個點對應一組有序數對 $(x,\,y)$，稱為它的\textbf{坐標（coordinates）}。於是「點」可以用「數」計算，「圖形」可以用「方程式」描述。本章都建立在這個對應上。

\subsection*{\sffamily B.\ 兩點距離與中點}

兩點 $A(x_1,\,y_1)$、$B(x_2,\,y_2)$ 之間的距離，是把水平差 $x_2-x_1$ 與鉛直差 $y_2-y_1$ 當作直角三角形的兩股，用\textbf{畢氏定理}求斜邊：
$$\boxed{\;\overline{AB}=\sqrt{(x_2-x_1)^2+(y_2-y_1)^2}\;}$$
$A$、$B$ 的\textbf{中點（midpoint）}$M$ 的坐標，是兩端坐標各取平均：
$$\boxed{\;M=\left(\frac{x_1+x_2}{2},\ \frac{y_1+y_2}{2}\right)\;}$$

\textbf{注意：}距離公式裡的差是平方，因此 $x_2-x_1$ 與 $x_1-x_2$ 結果相同，不必在意誰減誰；但稍後計算斜率時，分子分母的相減順序須一致（見 4-2）。中點是「平均」而非「相減」：$\left(\dfrac{x_2-x_1}{2},\dfrac{y_2-y_1}{2}\right)$ 不是中點，而是「半個位移量」。

\subsection*{\sffamily C.\ 對稱：把一個點照鏡子}

把一個點 $P(a,\,b)$ 對某條軸或某個點做\textbf{對稱（symmetry）}，相當於以那條軸或那個點為鏡子，求鏡中的像。其共同規則是「對稱軸是原點與像點連線的\textbf{垂直平分線}」，由此可推出下列四個常用結果：

\begin{center}
\renewcommand{\arraystretch}{1.35}
\begin{tabular}{p{0.26\linewidth} p{0.26\linewidth} p{0.30\linewidth}}
\hline
\textbf{對稱於} & \textbf{像點坐標} & \textbf{規則} \\
\hline
$x$ 軸 & $(a,\,-b)$ & $y$ 變號 \\
$y$ 軸 & $(-a,\,b)$ & $x$ 變號 \\
原點 & $(-a,\,-b)$ & $x,\,y$ 都變號 \\
直線 $y=x$ & $(b,\,a)$ & $x,\,y$ 對調 \\
\hline
\end{tabular}
\end{center}

\fig{d14-symmetry}{0.66}{圖 4-1：一點 $P$ 對 $x$ 軸、$y$ 軸、原點、直線 $y=x$ 的對稱點。}

\begin{補充框}{為什麼對 $y=x$ 對稱是「坐標對調」？}
直線 $y=x$ 上的點都滿足「橫坐標 $=$ 縱坐標」，它扮演把 $x$、$y$ 兩個角色互換的鏡子。要驗證 $P(a,b)$ 與 $(b,a)$ 對 $y=x$ 對稱，可檢查兩件事：
\begin{itemize}\setlength{\itemsep}{1pt}
\item 兩點的\textbf{中點}$\left(\frac{a+b}{2},\frac{a+b}{2}\right)$ 落在 $y=x$ 上（橫縱相等）；
\item 兩點\textbf{連線的斜率}$\dfrac{a-b}{b-a}=-1$，與 $y=x$ 的斜率 $1$ 乘積為 $-1$，兩者垂直。
\end{itemize}
「中點在軸上」加上「連線垂直於軸」，即表示該軸是連線的垂直平分線，也就是對稱。這也是「函數與它的反函數圖形對 $y=x$ 對稱」的原因：把 $(x,y)$ 換成 $(y,x)$ 正是反函數所做的事。
\end{補充框}

判斷一條\textbf{曲線}是否對某軸對稱，是看「整條曲線」而非單一點。例如把 $y=x^2$ 中的 $x$ 換成 $-x$，得 $y=(-x)^2=x^2$，方程式不變，因此 $y=x^2$ 對 $y$ 軸對稱。這是「拋物線左右對稱」的代數說法。

\section*{\sffamily 4-2\quad 直線方程式}

一條直線最關鍵的數字是它的\textbf{斜率}，即傾斜程度。先把斜率講清楚，後面的點斜式、平行與垂直、距離公式才有依據。

\subsection*{\sffamily A.\ 斜率：傾斜程度的數字化}

在直線上每往右走一單位（$x$ 增加），$y$ 上升或下降的量，即「\textbf{鉛直變化 $\div$ 水平變化}」這個比值，就是斜率。它不隨所取的兩點而改變，這正是「直線」的特徵。

\textbf{斜率（slope）}：直線上相異兩點 $A(x_1,y_1)$、$B(x_2,y_2)$（$x_1\ne x_2$），其斜率
$$\boxed{\;m=\frac{y_2-y_1}{x_2-x_1}=\frac{\Delta y}{\Delta x}\;}$$
\begin{itemize}\setlength{\itemsep}{1pt}
\item $m>0$：直線由左下往右上（遞增）；$m<0$：由左上往右下（遞減）。
\item $m=0$：水平線 $y=k$。鉛直線 $x=h$ 的\textbf{斜率不存在}（$\Delta x=0$，不能除以零）。
\item $|m|$ 越大越陡：同樣往右一單位，$y$ 變化越多。
\end{itemize}

\begin{補充框}{斜率為什麼能代表「傾斜程度」？$|m|$ 越大為什麼越陡？}
斜率量的是「\textbf{每往右走一單位，往上爬多少}」。固定往右走 $1$ 單位（$\Delta x=1$），則 $\Delta y=m$：$m=1$ 對應傾斜 $45^\circ$；$m=2$ 往右 $1$、往上 $2$，比 $45^\circ$ 更陡；$m=\frac12$ 往右 $1$、只往上 $\frac12$，較平緩。因此 $|m|$ 越大，同樣前進一步爬得越高，直線越接近鉛直。取絕對值是因為「陡不陡」只看傾斜程度、不看方向：$m=2$ 與 $m=-2$ 一樣陡，只是一個上升、一個下降。

\textbf{為什麼與選哪兩點無關？}在同一條直線上任取兩點作直角三角形，再取另外兩點作另一個直角三角形，因為四點都在同一條直線上，這兩個直角三角形\textbf{相似}，對應邊成比例，故 $\dfrac{\Delta y}{\Delta x}=\dfrac{\Delta y'}{\Delta x'}$。這個「比值固定不變」正是直線的定義性特徵；曲線沒有此性質（不同段陡度不同）。可類比為公路坡度牌「坡度 10\%」：每水平前進 $100$ 公尺上升 $10$ 公尺，即斜率 $0.1$。

當直線接近鉛直時 $\Delta x\to0$，比值沒有有限的數值，這就是\textbf{鉛直線斜率「不存在」}的原因。斜率其實是傾斜角 $\theta$（直線與 $x$ 軸正向的夾角）的正切 $m=\tan\theta$，此關聯在三角比一章正式處理。

\textbf{注意：}斜率與角度不是同一個數。斜率 $2$ 不代表「傾斜 $2$ 度」，而是 $\tan\theta=2$，對應 $\theta\approx63.4^\circ$；斜率加倍，角度並不加倍（$\tan$ 不是線性的）。另外，「斜率不存在」與「斜率為 $0$」不同：鉛直線是斜率不存在，水平線才是斜率 $=0$。
\end{補充框}

\subsection*{\sffamily B.\ 直線方程式的三種寫法}

已知斜率與一個點，整條直線就確定了。由斜率定義 $\dfrac{y-y_1}{x-x_1}=m$ 移項，即得最常用的\textbf{點斜式}：

\textbf{直線方程式的三式}（設斜率為 $m$）：
\begin{itemize}\setlength{\itemsep}{1pt}
\item \textbf{點斜式（point-slope form）}：過點 $(x_1,y_1)$，$\;y-y_1=m(x-x_1)$。
\item \textbf{斜截式（slope-intercept form）}：$y$ 截距為 $b$，$\;y=mx+b$。
\item \textbf{一般式（general form）}：$ax+by+c=0$（$a,b$ 不同時為 $0$）。
\end{itemize}
$$\boxed{\;y-y_1=m(x-x_1)\;}\qquad\boxed{\;y=mx+b\;}\qquad\boxed{\;ax+by+c=0\;}$$
一般式的優點是\textbf{鉛直線 $x=h$ 也能寫}（取 $b=0$）；當 $b\ne0$ 時其斜率為 $-\dfrac{a}{b}$。這三式是同一條直線的不同表示，可自由互換：有「點與斜率」用點斜式，要看截距用斜截式，要統一格式（尤其含鉛直線時）用一般式。

\subsection*{\sffamily C.\ 平移}

把直線（或任何圖形）\textbf{平移（translation）}：向右移 $h$、向上移 $k$，相當於把方程式中的 $x$ 換成 $x-h$、$y$ 換成 $y-k$。圖形 $f(x,y)=0$ 平移後的新方程式為
$$\boxed{\;f(x-h,\ y-k)=0\;}$$
其理由是：新圖上的點 $(x,y)$，其「平移前的位置」是 $(x-h,\,y-k)$，須以平移前的坐標滿足原方程式。直線平移後斜率不變、仍與原線平行，平移只改變截距。

\subsection*{\sffamily D.\ 平行線與垂直線}

設兩直線斜率分別為 $m_1,\,m_2$（皆存在）：
$$\boxed{\;\text{平行（parallel）}:\ m_1=m_2\;}\qquad\boxed{\;\text{垂直（perpendicular）}:\ m_1\cdot m_2=-1\;}$$
平行時斜率相等（且截距不同，否則重合）；垂直時兩斜率乘積為 $-1$，即互為\textbf{負倒數}。

\begin{補充框}{平行為什麼斜率相等？垂直為什麼乘積是 $-1$？}
\textbf{平行}：斜率是傾斜角的正切 $m=\tan\theta$。兩線平行表示它們與 $x$ 軸的傾斜角相同（$\theta_1=\theta_2$），所以 $\tan\theta_1=\tan\theta_2$，即 $m_1=m_2$。簡言之，平行 $=$ 同傾斜角 $=$ 同斜率。

\textbf{垂直}：取斜率 $m_1$ 的方向「往右 $1$、往上 $m_1$」。把這個方向整個\textbf{旋轉 $90^\circ$}，「往右 $1$、往上 $m_1$」會變成「往上 $1$、往左 $m_1$」，也就是「往右 $-m_1$、往上 $1$」。這條垂直線的斜率因此是
$$m_2=\frac{\text{上升}}{\text{前進}}=\frac{1}{-m_1}=-\frac{1}{m_1}\;\Longrightarrow\;m_1 m_2=-1.$$
那個 $-1$ 來自旋轉 $90^\circ$：負號來自方向反轉，倒數來自橫與縱對調。

\textbf{注意：}「垂直線斜率乘積是 $-1$」有一個前提——\textbf{兩條線斜率都存在}。水平線（$m=0$）與鉛直線（斜率不存在）也互相垂直，但不能寫成 $0\times(\text{不存在})=-1$。這是此判定的唯一例外。
\end{補充框}

特例：水平線（$m=0$）與鉛直線（斜率不存在）互相垂直。

\subsection*{\sffamily E.\ 點到直線的距離}

從一點 $P$ 向直線 $L$ 作\textbf{垂線}，垂足與 $P$ 的距離就是「點到直線的距離」，它是 $P$ 到 $L$ 上所有點之中\textbf{最短}的一段。其公式為：

\textbf{點到直線距離公式}：點 $P(x_0,\,y_0)$ 到直線 $L:\ ax+by+c=0$ 的距離
$$\boxed{\;d=\frac{|ax_0+by_0+c|}{\sqrt{a^2+b^2}}\;}$$
用法：把點代入直線左式取絕對值，再除以「$x,y$ 係數的平方和開根號」。

\fig{d14-point-line-distance}{0.56}{圖 4-2：點 $P(x_0,y_0)$ 到直線 $L:\,3x+4y-12=0$ 的距離 $d$ 是 $P$ 到垂足 $Q$ 的長度，即 $P$ 到 $L$ 上各點的最短距離。}

\textbf{注意：}套公式前直線一定要先寫成 $ax+by+c=0$（等號右邊為 $0$）。若直線是 $y=2x+3$，須先整理成 $2x-y+3=0$，才能讀出 $a=2,\,b=-1,\,c=3$；直接拿 $y=mx+b$ 套公式會出錯。分母是 $\sqrt{a^2+b^2}$，不是 $a^2+b^2$，也不含 $c$。

兩條平行線 $ax+by+c_1=0$、$ax+by+c_2=0$（係數 $a,b$ 已對齊）之間的距離為
$$\boxed{\;d=\frac{|c_1-c_2|}{\sqrt{a^2+b^2}}\;}$$
其理由是：平行線間的距離等於「任取一條線上一點，到另一條線的距離」。使用此式前，兩條線的 $a,\,b$ 必須先化成完全相同（例如 $6x+8y+1=0$ 要先約成 $3x+4y+\tfrac12=0$，才能與 $3x+4y+15=0$ 比較）。

\begin{延伸框}{點到直線距離公式是怎麼推出來的}
分子的 $ax_0+by_0+c$（把點代入直線左式）量的是「$P$ 偏離界線多少」，但它是用「斜的單位」量的；除以 $\sqrt{a^2+b^2}$ 就是把它換算回真正的長度單位。完整推導如下。

\textbf{關鍵物件——法向量。}直線 $ax+by+c=0$ 有一個天生的\textbf{法向量（normal vector）}$\vec n=(a,b)$，它\textbf{垂直}於這條直線。（簡單的驗證：直線上任兩點 $P_1,P_2$ 都滿足直線方程式，兩式相減得 $a(x_1-x_2)+b(y_1-y_2)=0$，這正是 $\vec n\cdot\overrightarrow{P_2P_1}=0$，故 $\vec n$ 垂直於直線方向。）點到直線的最短路徑，就是沿 $\vec n$ 方向走過去。

\textbf{推導。}設從 $P(x_0,y_0)$ 沿法向量方向走到垂足 $Q$，則
$$Q=(x_0+ta,\ y_0+tb),$$
其中 $t$ 是「走了幾個 $\vec n$ 的長度」的倍數。因為一個 $\vec n=(a,b)$ 的長度是 $\sqrt{a^2+b^2}$，故
\begin{equation}
d=\overline{PQ}=|t|\cdot\sqrt{a^2+b^2}.\tag{$*$}
\end{equation}
$Q$ 在直線上，代入 $ax+by+c=0$：
$$a(x_0+ta)+b(y_0+tb)+c=0\ \Longrightarrow\ (ax_0+by_0+c)+t(a^2+b^2)=0,$$
解得 $t=-\dfrac{ax_0+by_0+c}{a^2+b^2}$。代回 $(*)$：
$$d=|t|\sqrt{a^2+b^2}=\frac{|ax_0+by_0+c|}{a^2+b^2}\cdot\sqrt{a^2+b^2}=\frac{|ax_0+by_0+c|}{\sqrt{a^2+b^2}}.$$
因此 $\sqrt{a^2+b^2}$ 的真身就是法向量 $\vec n=(a,b)$ 的長度；分母的作用，是把「以法向量為尺」量出的偏離量，換算成真實距離。代入值的正負號還能判斷點在界線的哪一側，這也是二元一次不等式判半平面的依據。
\end{延伸框}

一條直線 $ax+by+c=0$ 把平面分成兩個\textbf{半平面（half-plane）}，不等式 $ax+by+c>0$ 與 $ax+by+c<0$ 各代表其中一邊。要畫出某個不等式的區域，作法是：先畫界線（取 $>$、$<$ 畫虛線表不含界線，取 $\ge$、$\le$ 畫實線含界線），再取一個不在線上的測試點（原點 $(0,0)$ 最方便，若它不在線上），代入左式判斷正負，使不等式成立的一側即為解區域。

\section*{\sffamily 4-3\quad 圓方程式}

「到一定點等距的所有點」構成圓。把這句話用兩點距離公式寫下來，就是圓的方程式。

\subsection*{\sffamily A.\ 圓的標準式}

圓心 $(h,\,k)$、半徑 $r$ 的圓，是所有「到圓心 $(h,k)$ 的距離等於 $r$」的點 $(x,y)$ 所成的集合。把這個條件 $\sqrt{(x-h)^2+(y-k)^2}=r$ 兩邊平方：
$$\boxed{\;(x-h)^2+(y-k)^2=r^2\;}$$
特例：圓心在原點時為 $x^2+y^2=r^2$。

\subsection*{\sffamily B.\ 一般式與配方還原圓心半徑}

把標準式展開，會得到 $x^2+y^2+Dx+Ey+F=0$ 的形式，稱為圓的\textbf{一般式}。反過來，看到一般式時須\textbf{配方（completing the square）}才能讀出圓心與半徑：
$$x^2+y^2+Dx+Ey+F=0\ \Longrightarrow\ \text{圓心}\left(-\tfrac D2,\,-\tfrac E2\right),\quad r=\sqrt{\tfrac{D^2+E^2}{4}-F}.$$
\textbf{重要前提}：根號內必須大於 $0$ 才是真正的圓；等於 $0$ 退化成一點，小於 $0$ 則無圖形。

\begin{補充框}{配方完別漏看右邊正負：圓的一般式退化三情況}
形如 $x^2+y^2+Dx+Ey+F=0$ 的式子\textbf{不一定}是圓。對 $x$、$y$ 各自配方後，會整理成
$$(x-h)^2+(y-k)^2=k_0,\qquad k_0=\frac{D^2+E^2}{4}-F,$$
其中左邊是兩個平方和、恆 $\ge0$。是否為圓、是什麼圖形，完全取決於右邊常數 $k_0$ 的正負，分為\textbf{三種情況}：
\begin{center}
\renewcommand{\arraystretch}{1.4}
\begin{tabular}{p{0.16\linewidth} p{0.30\linewidth} p{0.40\linewidth}}
\hline
\textbf{右邊 $k_0$} & \textbf{圖形} & \textbf{原因} \\
\hline
$k_0>0$ & \textbf{圓}，半徑 $r=\sqrt{k_0}$ & 「到圓心 $(h,k)$ 距離的平方 $=$ 正數」有一整圈解 \\
$k_0=0$ & \textbf{退化成一點} $(h,k)$ & 只有圓心本身使「距離平方 $=0$」，圓縮成一個點 \\
$k_0<0$ & \textbf{空集合}（無圖形） & 兩平方和恆 $\ge0$，不可能等於負數，沒有任何實點 \\
\hline
\end{tabular}
\end{center}
換言之，同一個一般式「照常配方」的步驟完全相同，但右邊常數的正負一變，結論就從「圓」變成「一個點」甚至「沒有圖形」。配方完務必檢查右邊 $k_0=\dfrac{D^2+E^2}{4}-F$ 是正、是零、還是負，再下結論。
\end{補充框}

\fig{d14-circle-degenerate}{0.96}{圖 4-4：圓的一般式配方後 $(x-h)^2+(y-k)^2=k_0$，同一圓心 $(h,k)$ 下右邊常數 $k_0$ 的三種情況——$k_0>0$ 是半徑 $\sqrt{k_0}$ 的圓，$k_0=0$ 退化成圓心一點，$k_0<0$ 沒有任何實點（空集合）。}

\textbf{注意：}配方時「加了多少要減回多少」，例如 $x^2-6x$ 要補上 $\big(\tfrac{-6}{2}\big)^2=9$ 湊成完全平方，就須在同一式減 $9$。此外，$x^2$ 與 $y^2$ 的係數須相等且為 $1$ 才是圓：看到 $2x^2+2y^2+\dots=0$，要先整條除以 $2$ 再配方；若係數不等（如 $x^2+4y^2=4$），那是橢圓而非圓。\textbf{最容易出錯的一點}是配方完只急著套圓心半徑、忘了看右邊 $k_0$ 的正負：題目常故意把 $F$ 調大，使「看似是圓的式子」其實是一個點（$k_0=0$）或無圖形（$k_0<0$），不檢查就會落入陷阱。

由給定條件求圓，常見三型：給圓心與半徑時，直接代標準式；給直徑兩端點 $A$、$B$ 時，圓心為 $AB$ 中點、半徑為 $\tfrac12\overline{AB}$；給圓上三點時，設一般式 $x^2+y^2+Dx+Ey+F=0$，三點各代入一次得三元一次聯立方程組，解出 $D,E,F$。

\section*{\sffamily 4-4\quad 直線與圓的關係}

一條直線與一個圓有三種位置關係：完全錯過（\textbf{相離}）、剛好相切（\textbf{相切}）、割進圓內（\textbf{相交}）。要判斷是哪一種，不必解聯立方程式，只要比較一個距離。

\subsection*{\sffamily A.\ 用「圓心到直線的距離」判定}

設圓半徑為 $r$，圓心到直線 $L$ 的距離為 $d$（以點到直線距離公式計算），則：
$$\boxed{\;d>r\Rightarrow\text{相離（0 交點）}\qquad d=r\Rightarrow\text{相切（1 交點）}\qquad d<r\Rightarrow\text{相交（2 交點）}\;}$$

\fig{d14-line-circle}{0.96}{圖 4-3：直線與圓的三種位置關係。比較圓心到直線的距離 $d$ 與半徑 $r$：$d>r$ 相離、$d=r$ 相切、$d<r$ 相交。}

\begin{補充框}{為什麼比「圓心到直線距離 $d$」與「半徑 $r$」就夠了？}
從圓心 $O$ 向直線 $L$ 作垂線，垂足 $H$ 是 $L$ 上離 $O$ 最近的點，$\overline{OH}=d$。沿這條垂線方向看，圓上離 $L$ 最近的點是垂線與圓的交點，它離 $L$ 的距離是 $d-r$（圓心離 $L$ 為 $d$，再靠近一個半徑 $r$）。於是：
\begin{itemize}\setlength{\itemsep}{1pt}
\item $d>r$：圓上最靠近 $L$ 的點離 $L$ 仍有 $d-r>0$，整個圓都搆不到 $L$，沒有交點（相離）。
\item $d=r$：圓恰好碰到 $L$ 一點（即垂足），這一點同時在圓上與線上，恰有一個交點（相切），且此時半徑 $OH$ 與切線垂直。
\item $d<r$：$L$ 從圓內穿過，割出兩個交點（一進一出）。
\end{itemize}
關鍵在於 $d$ 是圓心到直線的\textbf{最短距離}，因此拿它與半徑比較，就能決定圓有沒有任何部分搆得到直線；換成任何斜著量的距離都不行，只有最短的那個才能當判準。

這個幾何判定與「把直線代入圓、看判別式 $\Delta$」的代數作法完全等價：可證明 $\Delta$ 的符號與 $r^2-d^2$ 的符號永遠一致，所以 $d<r\Leftrightarrow\Delta>0$（兩交點）、$d=r\Leftrightarrow\Delta=0$（相切）、$d>r\Leftrightarrow\Delta<0$（相離）。兩者是同一件事的兩種說法：比 $d$ 與 $r$ 的幾何法快、又可直接用於設相切條件求參數；代入圓解判別式的代數法較慢，但能順便求出交點坐標。
\end{補充框}

\subsection*{\sffamily B.\ 圓的切線}

求圓的切線，依「切點是否已知」分兩種情況：
\begin{itemize}\setlength{\itemsep}{2pt}
\item \textbf{已知切點 $P(x_1,y_1)$ 在圓上}：切線與半徑垂直，可直接用切線公式，作法是把圓方程式中「一半的變數」換成切點坐標（降一次）：
\begin{itemize}\setlength{\itemsep}{1pt}
\item 圓心在原點 $x^2+y^2=r^2$：$\;\boxed{\;x_1 x+y_1 y=r^2\;}$
\item 圓心不在原點 $(x-h)^2+(y-k)^2=r^2$：$\;\boxed{\;(x_1-h)(x-h)+(y_1-k)(y-k)=r^2\;}$
\end{itemize}
（把 $(x-h)^2=(x-h)(x-h)$ 的其中一個因式換成 $(x_1-h)$，$y$ 同理；原點式是它在 $h=k=0$ 的特例。）
\item \textbf{已知切線斜率，或切線過某外部點}：設切線方程式，再用「圓心到它的距離 $=r$」求出未知參數。此法對任何圓（不論圓心在哪）都適用。
\end{itemize}
其依據是 4-4 A 的判定：相切等價於 $d=r$，因此「某直線與圓相切，求其中未知數」這型，標準作法就是令 $d=r$ 解方程式。此法比用判別式 $\Delta=0$ 快，因為不必先把直線代入圓展開。

\textbf{注意：}過\textbf{圓外}一點有\textbf{兩條}切線、過\textbf{圓上}一點有\textbf{一條}、圓\textbf{內}一點\textbf{沒有}切線；求切線前先判斷該點在圓內、圓上或圓外。又：用「設 $y=mx+k$」求切線時會漏掉鉛直切線（斜率不存在的那條），須另外檢查鉛直線 $x=$ 常數是否也相切。

\section*{\sffamily 4-5\quad 本章重點整理}

\begin{補充框}{一頁複習（公式與關鍵結論）}
\begin{itemize}\setlength{\itemsep}{2pt}
\item \textbf{距離與中點}：$\overline{AB}=\sqrt{(x_2-x_1)^2+(y_2-y_1)^2}$；中點取兩端坐標的平均。
\item \textbf{對稱}：對 $x$ 軸 $(a,-b)$、$y$ 軸 $(-a,b)$、原點 $(-a,-b)$、$y=x$ 對調為 $(b,a)$。
\item \textbf{斜率}：$m=\dfrac{\Delta y}{\Delta x}=\tan\theta$；$|m|$ 越大越陡；鉛直線斜率不存在（不等於 $0$）。
\item \textbf{直線三式}：點斜式 $y-y_1=m(x-x_1)$、斜截式 $y=mx+b$、一般式 $ax+by+c=0$。
\item \textbf{平行與垂直}：$m_1=m_2$ 平行；$m_1 m_2=-1$ 垂直（互為負倒數，前提為兩斜率都存在）。
\item \textbf{點到直線距離}：$d=\dfrac{|ax_0+by_0+c|}{\sqrt{a^2+b^2}}$；平行線間距 $\dfrac{|c_1-c_2|}{\sqrt{a^2+b^2}}$（係數先對齊）。
\item \textbf{圓}：標準式 $(x-h)^2+(y-k)^2=r^2$；一般式 $x^2+y^2+Dx+Ey+F=0$ 配方還原圓心與半徑。配方完看右邊 $\tfrac{D^2+E^2}{4}-F$：$>0$ 才是圓、$=0$ 退化成一點、$<0$ 是空集合。
\item \textbf{直線與圓}：比 $d$ 與 $r$——$d>r$ 相離、$d=r$ 相切、$d<r$ 相交；圓上一點切線：原點圓 $x_1 x+y_1 y=r^2$、一般圓 $(x_1-h)(x-h)+(y_1-k)(y-k)=r^2$，外部點或已知斜率用 $d=r$ 求（任何圓皆可）。
\end{itemize}
\end{補充框}

\end{document}
